\newcommand{\CircuitoCorrentiMaglia}[1]{%
\begin{figure}[h]
    \centering
    \begin{circuitikz}[scale=1.2, every node/.style={font=\small}]
        % Parte fissa (resistori e collegamenti)
        \draw (0.75,2) to[american resistor, l_={$R_2$}] (0.75,4);
        \draw (2.25,2) to[american resistor, l={$R_3$}] (4,2);
        \draw (4,4) to[american resistor, l={$R_4$}] (4,2);
        \draw (0.75,0) to[american resistor, l={$R_1$}] (4,0);
        \draw (0.75,2) -- (0.75,0);
        \draw (6,2) to[american resistor, l_={$R_5$}] (6,0);
        \draw (4,0) -- (6,0);
        \draw (4,4) -- (6,4);

        % Parte variabile: generatori
        \ifcase#1
            % caso 0 (non usato)
        \or % #1 = 1: soli generatori di tensione, senza frecce
            \draw (0.75,4) to[american voltage source, l={$E_2$}] (4,4);
            \draw (2.25,2) to[american voltage source, l={$E_1$}] (0.75,2);
            \draw (4,0) to[american voltage source, l={$E_3$}] (4,2);
            \draw (6,4) to[american voltage source, l_={$E_4$}] (6,2);
        \or % #1 = 2: soli generatori di tensione, con frecce
            \draw (0.75,4) to[american voltage source, l={$E_2$}] (4,4);
            \draw (2.25,2) to[american voltage source, l={$E_1$}] (0.75,2);
            \draw (4,0) to[american voltage source, l={$E_3$}] (4,2);
            \draw (6,4) to[american voltage source, l_={$E_4$}] (6,2);
            % Frecce correnti di maglia
            \draw[->, thick, red] (1.7,1) arc (120:60:1.2) node[midway, above]{$I_1$};
            \draw[->, thick, blue] (1.7,3) arc (120:60:1.2) node[midway, above]{$I_2$};
            \draw[->, thick, magenta] (5,1) arc (240:110:1) node[midway, right]{$I_3$};
        \or % #1 = 3: due generatori di corrente (J1,J2), senza frecce
            \draw (0.75,4) to[american current source, l={$J_1$}] (4,4);
            \draw (2.25,2) to[american voltage source, l={$E_1$}] (0.75,2);
            \draw (4,0) to[american current source, l={$J_2$}] (4,2);
            \draw (6,4) to[american voltage source, l_={$E_4$}] (6,2);
        \or % #1 = 4: due generatori di corrente (J1,J2), con frecce
            \draw (0.75,4) to[american current source, l={$J_1$}] (4,4);
            \draw (2.25,2) to[american voltage source, l={$E_1$}] (0.75,2);
            \draw (4,0) to[american current source, l={$J_2$}] (4,2);
            \draw (6,4) to[american voltage source, l_={$E_4$}] (6,2);
            % Frecce correnti di maglia
            \draw[->, thick, red] (1.7,1) arc (120:60:1.2) node[midway, above]{$I_1$};
            \draw[->, thick, blue] (1.7,3) arc (120:60:1.2) node[midway, above]{$I_2$};
            \draw[->, thick, magenta] (5,1) arc (240:110:1) node[midway, right]{$I_3$};
        \fi

        % Nodi (comuni a tutte le figure)
        \node[circ] at (4,4){};
        \node at (4,4) [above] {A};
        \node[circ] at (4,0){};
        \node at (4,0) [below] {B};
        \node[circ] at (0.75,2){};
        \node at (0.75,2) [left] {C};
        \node[circ] at (4,2){};
        \node at (4,2) [right] {D};
    \end{circuitikz}
    \caption{%
        \ifcase#1
        \or Circuito Esempio
        \or Circuito con tre maglie elementari. Le correnti di maglia \(I_1\), \(I_2\) e \(I_3\) sono in senso orario.
        \or Circuito con tre maglie e due generatori di corrente
        \or Circuito con tre maglie e due generatori di corrente. Le correnti di maglia sono in senso orario.
        \fi
    }
    \label{fig:circuito_#1}
\end{figure}
}


\newcommand{\CircuitoPotenzialeNodi}[1]{%
\begin{figure}[h]
    \centering
    \begin{circuitikz}[scale=1.2, every node/.style={font=\small}]
        % Parte fissa (resistori e collegamenti)
        \draw (0.75,2) to[american resistor, l_={$R_2$}] (0.75,4);
        \draw (2.25,2) to[american resistor, l={$R_3$}] (4,2);
        \draw (4,4) to[american resistor, l={$R_4$}] (4,2);
        \draw (0.75,0) to[american resistor, l={$R_1$}] (4,0);
        \draw (0.75,2) -- (0.75,0);
        \draw (6,2) to[american resistor, l_={$R_5$}] (6,0);
        \draw (4,0) -- (6,0);
        \draw (4,4) -- (6,4);
        
        \ifcase#1
            % caso 0 non usato
        \or % #1 = 1
            \draw (0.75,4) to[american current source, l={$J_1$}] (4,4);
            \draw (2.25,2) to[american current source, l={$J_2$}] (0.75,2);
            \draw (4,0) to[american current source, l={$J_3$}] (4,2);
            \draw (6,4) to[american current source, l_={$J_4$}] (6,2);
            % senza ground
        \or % #1 = 2
            \draw (0.75,4) to[american current source, l={$J_1$}] (4,4);
            \draw (2.25,2) to[american current source, l={$J_2$}] (0.75,2);
            \draw (4,0) to[american current source, l={$J_3$}] (4,2);
            \draw (6,4) to[american current source, l_={$J_4$}] (6,2);
            \node[ground] at (4,0) {};
        \or % #1 = 3
            \draw (0.75,4) to[american current source, l={$J_1$}] (4,4);
            \draw (2.25,2) to[american voltage source, l={$E_1$}] (0.75,2);
            \draw (4,0) to[american current source, l={$J_2$}] (4,2);
            \draw (6,4) to[american voltage source, l_={$E_4$}] (6,2);
            % senza ground
        \or % #1 = 4
            \draw (0.75,4) to[american current source, l={$J_1$}] (4,4);
            \draw (2.25,2) to[american voltage source, l={$E_1$}] (0.75,2);
            \draw (4,0) to[american current source, l={$J_2$}] (4,2);
            \draw (6,4) to[american voltage source, l_={$E_4$}] (6,2);
            \node[ground] at (4,0) {};
        \fi
        
        % Nodi (etichette)
        \node[circ] at (4,4){}; \node at (4,4) [above] {A};
        \node[circ] at (0.75,2){}; \node at (0.75,2) [left] {C};
        \node[circ] at (4,2){}; \node at (4,2) [right] {D};
        \node[circ] at (4,0){}; \node at (4,0.2) [right] {B};
    \end{circuitikz}
    \caption{%
        \ifcase#1
        \or Circuito per il metodo dei potenziali ai nodi (solo generatori di corrente) senza GND
        \or Circuito per il metodo dei potenziali ai nodi (solo generatori di corrente) con GND
        \or Circuito per il metodo dei potenziali ai nodi (generatori di corrente e tensione) senza GND
        \or Circuito per il metodo dei potenziali ai nodi (generatori di corrente e tensione) con GND
        \fi
    }
    \label{fig:potenziale_#1}
\end{figure}
}





\usetikzlibrary{decorations.markings}

\usetikzlibrary{arrows.meta}

\newcommand{\ZoomNodo}[1]{%
\begin{center}

\begin{circuitikz}[scale=1.1, decoration={markings, mark=at position 0.5 with {\arrow{Stealth[scale=2]}}}]
    \ifstrequal{#1}{A}{%
        \node[circ] (A) at (0,0) {};
        \node[above right] at (A) {A};
        \draw[postaction={decorate}] (A) -- (1.5,0) node[midway, above] {$J_4$};
        \draw[postaction={decorate}] (-1.5,0) -- (A) node[midway, above] {$J_1$};
        \draw[postaction={decorate}] (A) -- (0,-1.5) node[midway, right] {$I_{R_4}$};
    }{}
    \ifstrequal{#1}{B}{%
        \node[ground] (B) at (0,0) {};
        \node[below] at (B) {B};
        \draw[postaction={decorate}] (B) -- (0,1.5) node[midway, right] {$I_{R1}$};
        \draw[postaction={decorate}] (B) -- (1.5,0) node[midway, above] {$I_{J2}$};
        \draw[postaction={decorate}] (B) -- (-1.5,0) node[midway, above] {$I_{E4,R5}$};
    }{}
    \ifstrequal{#1}{C}{%
        \node[circ] (C) at (0,0) {};
        \node[above left] at (C) {C};
        \draw[postaction={decorate}] (C) -- (0,-1.5) node[midway, right] {$I_{R_1}$};
        \draw[postaction={decorate}] (C) -- (0,1.5) node[midway, right] {$J_1$};
        \draw[postaction={decorate}] (1.5,0) -- (C) node[midway, above] {$J_2$};
    }{}
    \ifstrequal{#1}{D}{%
        \node[circ] (D) at (0,0) {};
        \node[above right] at (D) {D};
        \draw[postaction={decorate}] (0,-1) -- (D) node[midway, right] {$J_3$};
        \draw[postaction={decorate}] (D) -- (0,1.5) node[midway, right] {$I_{R_4}$};
        \draw[postaction={decorate}] (D) -- (-1.5,0) node[midway, above] {$J_2$};
    }{}
\end{circuitikz}
\end{center}
}

    